% Part: second-order-logic
% Chapter: metatheory
% Section: compactness

\documentclass[../../../include/open-logic-section]{subfiles}

\begin{document}

\olfileid{sol}{met}{lst}

\olsection{द्वितीय-क्रम तर्कशास्त्रात लोव्हेनहाइम--स्कोलेम प्रमेय लागू होत नाही}

\begin{explain}
(अधोगामी) लोव्हेनहाइम--स्कोलेम प्रमेयानुसार, अनंत
प्रतिमान असलेल्या !!{sentence}s च्या प्रत्येक संचाला
!!a{enumerable} प्रतिमानही असते. हे प्रमेयदेखील
संपूर्णता प्रमेयाचा परिणाम आहे: संपूर्णतेचा पुरावा
!!{sentence}s च्या कोणत्याही सुसंगत संचासाठी प्रतिमान
तयार करतो आणि ते प्रतिमान !!{enumerable} असते.
ऊर्ध्वगामी लोव्हेनहाइम--स्कोलेम प्रमेयही आहे.
त्यानुसार, !!{sentence}s च्या संचाला
!!a{denumerable} प्रतिमान असेल, तर त्याला
!!a{nonenumerable} प्रतिमानही असते. ही दोन्ही
प्रमेये द्वितीय-क्रम तर्कशास्त्रात लागू होत नाहीत.
\end{explain}


\begin{thm}
\ollabel{thm:sol-no-ls} लोव्हेनहाइम--स्कोलेम प्रमेय
द्वितीय-क्रम तर्कशास्त्रात लागू होत नाही:
काही !!{sentence}s ना अनंत प्रतिमाने असतात,
पण !!{enumerable} प्रतिमाने नसतात.
\end{thm}

\begin{proof}
हे आठवा:
\[
\fn{Count} \ident \lexists[z][\lexists[u][\lforall[X][((X(z) \land
      \lforall[x][(X(x) \lif X(u(x)))]) \lif \lforall[x][X(x)])]]]
\]
हे !!a{structure}~$\Struct{M}$ मध्ये सत्य असेल,
जर आणि फक्त जर $\Domain{M}$ !!{enumerable} असेल.
म्हणून $\lnot\fn{Count}$ हे~$\Struct{M}$ मध्ये
सत्य असेल, जर आणि फक्त जर $\Domain{M}$
!!{nonenumerable} असेल. अशा !!{structure}s
आहेत---कोणताही !!{nonenumerable} संच !!{domain}
म्हणून घ्या; उदाहरणार्थ, $\Pow{\Nat}$ किंवा~$\Real$.
म्हणून $\lnot \fn{Count}$ ला अनंत प्रतिमाने
आहेत, पण !!{enumerable} प्रतिमाने नाहीत.
\end{proof}

\begin{thm}
काही !!{sentence}s ना !!{denumerable} प्रतिमाने
असतात, पण !!{nonenumerable} प्रतिमाने नसतात.
\end{thm}

\begin{proof}
$\fn{Count} \land \fn{Inf}$ हे $\Nat$ मध्ये
सत्य असते; पण $\Domain{M}$ !!{nonenumerable}
असलेल्या कोणत्याही !!{structure}~$\Struct{M}$
मध्ये ते सत्य नसते.
\end{proof}


\end{document}
